\clearpage\SourceLeaf{450}%
\begin{center}\begingroup\fontsize{8}{9}\selectfont
\setlength{\tabcolsep}{4pt}
\begin{tabular}{r|r|p{170pt}|r}\hline
\(+\) south&\(+\) west&&No.\\\hline
\(-242773.3\)&\(+36827.3\)&Drochtersen&\((1)\)\\
\(-235431.6\)&\(+29726.7\)&Colmar&\((1)\)\\
\(-245551.2\)&\(+90763.0\)&Cappeln&\((1)\)\\
\(-245775.5\)&\(+87628.9\)&Midlum&\((1)\)\\
\(-245821.6\)&\(+24600.1\)&Neuendorf&\((1)\)\\
\(-246532.7\)&\(+67702.1\)&Oster Ilienworth&\((1)\)\\
\(-246892.5\)&\(+76614.2\)&Wanna&\((1)\)\\
\(-247388.5\)&\(+63676.7\)&Bülkau&\((1)\)\\
\(-247703.2\)&\(+36752.9\)&Krautsand&\((1)\)\\
\(-250722.6\)&\(+69097.6\)&Neuenkirchen&\((1)\)\\
\(-250784.2\)&\(+65963.9\)&Osterbruch&\((1)\)\\
\(-251343.0\)&\(+34554.1\)&Glückstadt small tower (Danish station)&\((1)\)\\
\(-251399.8\)&\(+30139.2\)&Herzhorn&\((1)\)\\
\(-251523.1\)&\(+11052.9\)&Barmstedt&\((1)\)\\
\(-251573.4\)&\(+34126.7\)&Glückstadt church tower&\((1)\)\\
\(-252146.1\)&\(+40748.4\)&Hammelvörden&\((1)\)\\
\(-253330.1\)&\(+78043.0\)&Lüdingworth&\((1)\)\\
\(-254110.4\)&\(+21378.6\)&Horst&\((1)\)\\
\(-254115.7\)&\(+68778.6\)&Otterndorf&\((1)\)\\
\(-255131.8\)&\(+28321.4\)&Süderau&\((1)\)\\
\(-255141.2\)&\(+62459.0\)&Belum&\((1)\)\\
\(-255213.6\)&\(+83123.0\)&Franzenburg&\((1)\)\\
\(-255856.6\)&\(+43050.3\)&Freiburg&\((1)\)\\
\hline\end{tabular}\endgroup\end{center}
\clearpage
\begin{center}\begingroup\fontsize{8}{9}\selectfont
\setlength{\tabcolsep}{4pt}
\begin{tabular}{r|r|p{170pt}|r}\hline
\(+\) south&\(+\) west&&No.\\\hline
\(-255908.9\)&\(+33800.4\)&Borsfleth&\((1)\)\\
\(-256062.6\)&\(+84523.8\)&Hartmann’s station 1 near Altenwalde&\((1)\)\\
\(-256062.7\)&\(+84104.8\)&Altenwalde&\((1)\)\\
\(-256069.1\)&\(+77120.7\)&Altenbruch spire 1.&\((1)\)\\
\(-256075.2\)&\(+77119.1\)&Altenbruch spire 2.&\((1)\)\\
\(-256894.2\)&\(+29816.6\)&Crempe&\((1)\)\\
\(-257722.7\)&\(+21032.0\)&Hohenfelde&\((1)\)\\
\(-257778.1\)&\(+35721.5\)&Wevelsfleth&\((1)\)\\
\(-257996.5\)&\(+80031.2\)&Groden&\((1)\)\\
\(-258998.8\)&\(+15610.9\)&Hörnerkirchen&\((1)\)\\
\(-259523.9\)&\(+26920.1\)&Neuenbrook&\((1)\)\\
\(-259779.2\)&\(+40370.0\)&Brockdorf&\((1)\)\\
\(-259854.2\)&\(+81878.1\)&Ritzebüttel gable rod 1&\((1)\)\\
\(-260086.9\)&\(+33836.4\)&Neuenkirchen&\((1)\)\\
\(-261494.4\)&\(+81196.7\)&Cuxhaven lighthouse&\((1)\)\\
\(-261725.0\)&\(+34726.2\)&Beienfleth&\((1)\)\\
\(-262754.7\)&\(+83558.8\)&Döse&\((1)\)\\
\(-263188.1\)&\(+45161.9\)&St. Margareth&\((1)\)\\
\(-266552.1\)&\(+37341\,8\)&Wilster&\((1)\)\\
\(-266569.396\)&\(+95074.242\)&Neuwerk lighthouse center&\((1)\)\\
\(-269380.6\)&\(+14702.3\)&Kellinghusen&\((1)\)\\
\(-270223.4\)&\(+61150.8\)&Marne&\((1)\)\\
\hline\end{tabular}\endgroup\end{center}

The following remarks explain the meaning of the coordinates.

For a general idea of their meaning, the first number can be taken to indicate how far the place in question lies south (with a \(+\) sign) or north (with a \(-\) sign) of the Göttingen Observatory. The second indicates how far it lies west (with \(+\)) or east (with \(-\)).

The curvature of the Earth’s surface has already been taken into account in such a way that, when these coordinates are plotted on a plane, the resulting image is \textit{conformal}, that is, similar in its smallest parts. My geodetic papers already give some details on this, and later papers will develop the subject more fully.

The exact origin of the coordinates at the observatory, incidentally, is the center of the axis of the Reichenbach meridian circle.

The unit chosen for the coordinates is the length that, to the best of the knowledge available in 1821, could be taken as the ten-millionth part of the quadrant of the Earth’s meridian, namely 443,307885 Paris lines. This differs slightly, though very little, from the so-called legal French metre. As is well known, the latter had been set at 443,296 Paris lines. Although further determinations of the ten-millionth part of the Earth’s meridian quadrant have been obtained in later years (since 1821), and indeed have always been distinctly greater than the legal metre just mentioned), I have nevertheless preferred to keep the unit I originally chose, since each individual number can easily be converted to any other unit once the ratio is known.

